\documentclass[11pt]{article}

\usepackage[utf8]{inputenc}
\usepackage[a4paper,margin=1in]{geometry}
\usepackage{graphicx}
\usepackage{amsmath}
\usepackage{amssymb}
\usepackage{amsthm}
\usepackage{booktabs}
\usepackage{url}
\usepackage[round]{natbib}
\usepackage[hidelinks]{hyperref}
\usepackage{microtype}
\setlength{\emergencystretch}{2em}

\theoremstyle{definition}
\newtheorem{definition}{Definition}
\newtheorem{example}{Example}

\theoremstyle{plain}
\newtheorem{theorem}{Theorem}
\newtheorem{proposition}[theorem]{Proposition}
\newtheorem{corollary}[theorem]{Corollary}
\newtheorem{lemma}[theorem]{Lemma}

\theoremstyle{remark}
\newtheorem{remark}{Remark}

\newcommand{\Auth}{\mathit{Auth}}
\newcommand{\RegionOK}{\mathit{RegionAllowed}}
\newcommand{\PurpOK}{\mathit{PurposeLimited}}
\newcommand{\MinData}{\mathit{MinData}}
\newcommand{\LegalOK}{\mathit{LegallyValid}}
\newcommand{\Act}{\mathit{Act}}
\newcommand{\SoD}{\mathit{SoD}}
\newcommand{\Tdom}{\mathbb{T}}
\newcommand{\Ldom}{\mathcal{L}}
\newcommand{\Pdom}{\mathcal{P}}

\title{Spatio-Temporal Purpose-Aware RBAC\\ for Federated Data Processing:\\ Model, Analysis, and Verification}
\author{Natallia Kokash}
\date{August 2026}

\begin{document}

\maketitle

\begin{abstract}
Federated data processing (FDP) networks execute workflows across organisations
that are subject to different legal regimes, on data whose permitted uses are
bounded by purpose, jurisdiction, and time. Classical role-based access control
cannot express these constraints, and the temporal and spatial extensions
proposed in the literature treat them in isolation. We present STP-RBAC, a
spatio-temporal purpose-aware access control model in which authorisation is a
predicate over a user, a role, an object, a declared purpose, a declared
operating jurisdiction, and a time instant. The model is given a temporal graph
semantics, and we establish four structural results: authorisation at a fixed
purpose, jurisdiction, and instant reduces to reachability in a classical RBAC
snapshot; authorisation is monotone under a policy relaxation order, from which
compliance verification inherits a downward-closure property that supports
incremental re-verification when participants join; temporal separation of duty
reduces to interval stabbing over the arrangement of activation endpoints; and
workflow feasibility reduces to time-respecting reachability, which is
polynomial for fixed role assignment and NP-complete once roles may be chosen
per step. We give an SMT encoding covering all four constraint dimensions, and
we distinguish policy checking, which is polynomial, from policy synthesis and
repair, which is not. We close with an evaluation design for the reference
implementation.
\end{abstract}

% ============================================================
\section{Introduction}
\label{sec:intro}

Federated data processing moves computation to data that cannot be centralised,
typically because the data are personal, commercially sensitive, or subject to
localisation requirements. In our earlier work we described an architecture for
compliance-aware FDP built on the Brane orchestration
framework~\citep{valkering2021brane}, together with a coordination platform
that elicits regulatory constraints from participants and helps translate them
into machine-actionable policies~\citep{kokash2026compliance}. That work
sketched a formal model, the Spatio-Temporal Purpose-Aware RBAC Graph, but left
its analysis open. The present paper develops that model into a complete theory
and specifies how to verify policies expressed in it.

The setting imposes four constraint dimensions that must hold simultaneously.
Access is granted for a bounded period, because ethics approvals, data use
agreements, and project phases all expire. Access depends on where the
requesting participant operates and where the data are governed and stored,
because adequacy decisions and localisation rules are asymmetric between
jurisdictions. Access is bound to a declared purpose, because purpose
limitation is a first-order obligation under the GDPR and comparable regimes.
Finally, only the objects that the declared purpose actually requires may be
touched, which is data minimisation. Existing extensions of RBAC address these
dimensions separately. TRBAC~\citep{bertino2001trbac} and
GTRBAC~\citep{joshi2005gtrbac} add time; GEO-RBAC~\citep{bertino2005georbac}
and LoT-RBAC~\citep{chandran2005lotrbac} add location; purpose-based access
control~\citep{byun2008purpose} adds intent. Combining them is not merely a
matter of conjoining predicates, because the dimensions interact: role
hierarchies can launder jurisdictional restrictions, and a policy that is safe
for each object in isolation can be unsafe for a workflow that touches several
objects in sequence.

\paragraph{Contributions.} This paper makes the following contributions.

\begin{enumerate}
  \item A unified model, STP-RBAC (Section~\ref{sec:model}), that adds users,
        user-role assignment, and a temporally scoped role hierarchy to the
        graph model of \citet{kokash2026compliance}. Users are necessary rather
        than cosmetic: separation of duty is a constraint on what one principal
        may hold, and it cannot be stated in a model whose vertices are only
        roles and objects.
  \item A corrected treatment of jurisdiction. We separate the regimes that
        \emph{govern} an object from the sites that \emph{store} it, and
        quantify universally over the former and existentially over the latter.
        The single existential quantifier used in earlier formulations admits a
        forum-shopping leak, described in Remark~\ref{rem:forum}.
  \item Four structural results (Section~\ref{sec:results}) with proofs:
        snapshot reduction, monotonicity under policy relaxation with its
        downward-closure corollary for compliance, an interval characterisation
        of temporal separation of duty, and a reduction of workflow feasibility
        to time-respecting reachability.
  \item An algorithmic treatment (Section~\ref{sec:algorithms}) with defensible
        complexity bounds, and an SMT encoding (Section~\ref{sec:smt}) that
        covers time, space, purpose, and minimisation, separating the
        polynomial checking problem from the hard synthesis problem.
  \item A verification architecture (Section~\ref{sec:architecture}) and an
        evaluation design (Section~\ref{sec:evaluation}) for the reference
        implementation.
\end{enumerate}

Section~\ref{sec:related} positions the work, and Section~\ref{sec:conclusion}
discusses limitations.

% ============================================================
\section{Setting and Requirements}
\label{sec:setting}

We summarise the deployment context assumed throughout. An FDP network consists
of a coordinator and a set of autonomous domains. Data providers hold local
datasets and run worker nodes that execute approved containerised steps inside
their own administrative boundary. Data consumers submit workflows to the
coordinator and receive only authorised outputs. Raw data never leave the
provider domain. Policies are enforced at two points: the coordinator, which
sees the workflow and the requesting principal but not the data, and the worker
node, which sees the request in full context. This split is architectural
rather than incidental, and Corollary~\ref{cor:split} shows that the model
respects it.

The following requirements drive the model.

\begin{description}
  \item[R1 (Bounded validity).] Every assignment carries an explicit validity
        window, and windows may be discontinuous, since project phases recur.
  \item[R2 (Asymmetric jurisdiction).] Whether a participant in jurisdiction
        $\ell$ may process data governed by jurisdiction $\ell'$ is a directed
        relation reflecting adequacy decisions and bilateral agreements.
  \item[R3 (Purpose binding).] Every request carries a declared purpose, and
        authorisation depends on it. Purposes cannot be compiled away, because
        they vary per request.
  \item[R4 (Minimisation).] A request may touch only objects that the declared
        purpose requires.
  \item[R5 (Workflow-level reasoning).] The unit of interest is often a
        multi-step workflow rather than a single access, so the model must
        support feasibility questions over sequences of requests.
\end{description}

% ============================================================
\section{The STP-RBAC Model}
\label{sec:model}

\subsection{Signature}

Fix finite sets $U$ of users, $R$ of roles, $O$ of data objects, $\Ldom$ of
jurisdictional labels, and $\Pdom$ of processing purposes, together with a
totally ordered time domain $\Tdom$. Let $\mathcal{I}$ denote the set of finite
unions of closed intervals over $\Tdom$; a periodic activation expression in the
sense of GTRBAC~\citep{joshi2005gtrbac} unfolds to an element of $\mathcal{I}$
over any bounded planning horizon, which is all we require.

\begin{definition}[STP-RBAC policy]
\label{def:policy}
An \emph{STP-RBAC policy} is a tuple
\[
  \Theta = \bigl(
    U,\, R,\, O,\, \Ldom,\, \Pdom,\;
    UA,\, PA,\, H,\;
    \tau_U,\, \tau_P,\, \tau_H,\;
    \rho_U,\, \rho_G,\, \rho_S,\, \Gamma,\, \Pi,\, \Delta,\, \SoD
  \bigr)
\]
where
\begin{itemize}
  \item $UA \subseteq U \times R$ is the user-role assignment and
        $\tau_U : UA \to \mathcal{I}$ its validity;
  \item $PA \subseteq R \times O$ is the permission assignment and
        $\tau_P : PA \to \mathcal{I}$ its validity;
  \item $H \subseteq R \times R$ is the role hierarchy, where $(r,r') \in H$
        means that $r$ is senior to $r'$, and $\tau_H : H \to \mathcal{I}$ is
        its validity;
  \item $\rho_U : U \to 2^{\Ldom} \setminus \{\emptyset\}$ gives the
        jurisdictions from which a user may operate;
  \item $\rho_G : O \to 2^{\Ldom} \setminus \{\emptyset\}$ gives the
        jurisdictions whose law \emph{governs} an object;
  \item $\rho_S : O \to 2^{\Ldom} \setminus \{\emptyset\}$ gives the
        jurisdictions in which an object is \emph{stored};
  \item $\Gamma \subseteq \Ldom \times \Ldom$ is the cross-jurisdiction access
        relation, where $(\ell, \ell') \in \Gamma$ means a participant
        operating in $\ell$ may process data attached to $\ell'$. $\Gamma$ is
        reflexive but not assumed symmetric or transitive;
  \item $\Pi : PA \to 2^{\Pdom}$ gives the purposes for which a role may
        process an object. Defining $\Pi$ on pairs rather than on roles alone
        captures object-specific restrictions, for instance an imaging dataset
        that a data manager may anonymise but not export;
  \item $\Delta : \Pdom \to 2^{O}$ gives the \emph{purpose-admissible object
        set}: the objects that purpose $p$ may legitimately require;
  \item $\SoD \subseteq R \times R$ is an irreflexive, symmetric set of
        mutually exclusive role pairs.
\end{itemize}
\end{definition}

\begin{remark}[Naming of $\Delta$]
\citet{kokash2026compliance} call $\Delta(p)$ the ``minimal required set'' of
objects for $p$, while testing $O' \subseteq \Delta(p)$. A minimal set under
a subset test is the wrong reading, since the test admits any subset including
the empty one. The predicate is an upper bound on what $p$ may touch, so we
rename it to purpose-admissible set. The formal content is unchanged; only the
name was misleading.
\end{remark}

\subsection{Requests and the temporal authorisation graph}

\begin{definition}[Access request]
An \emph{access request} is a tuple $q = (u, r, o, p, \ell, t)$ in
$U \times R \times O \times \Pdom \times \Ldom \times \Tdom$, read as: user $u$,
acting in role $r$ from jurisdiction $\ell$, requests object $o$ for purpose $p$
at time $t$.
\end{definition}

Making $\ell$ part of the request, rather than deriving it from the role, is a
deliberate change from \citet{kokash2026compliance}. It puts jurisdiction on
the same footing as purpose: both are declared per request, both are attested by
the requester, and both are checked rather than assumed. It also removes the
existential quantifier over the requester's possible locations, which is what
Remark~\ref{rem:forum} is about.

For an assignment $x$ with validity $\tau(x)$, write
$\chi_x(t) = 1$ if $t \in \tau(x)$ and $0$ otherwise. The
\emph{active hierarchy} at time $t$ is
$H_t = \{(r,r') \in H \mid \chi_{(r,r')}(t) = 1\}$, and $\succeq_t$ denotes the
reflexive-transitive closure of $H_t$. We require $H$ to be acyclic, so
$\succeq_t$ is a partial order for every $t$.

\begin{definition}[Effective assignment and effective purposes]
\label{def:effective}
The \emph{effective permission assignment} at time $t$ is
\[
  PA^{*}_t = \bigl\{ (r,o) \;\big|\;
    \exists r' \in R:\; r \succeq_t r',\;
    (r',o) \in PA,\; \chi_{(r',o)}(t) = 1 \bigr\},
\]
and the \emph{effective purpose set} of $(r,o) \in PA^{*}_t$ is
\[
  \Pi_t(r,o) = \bigcup \bigl\{ \Pi(r',o) \;\big|\;
    r \succeq_t r',\; (r',o) \in PA,\; \chi_{(r',o)}(t) = 1 \bigr\}.
\]
The \emph{activation closure} of a user is
\[
  \Act_t(u) = \bigl\{ r \in R \;\big|\;
    \exists r' :\; (u,r') \in UA,\; \chi_{(u,r')}(t) = 1,\; r' \succeq_t r
  \bigr\}.
\]
\end{definition}

Permissions flow upward through the hierarchy and role membership flows
downward, which is the standard RBAC reading. What does \emph{not} flow is
jurisdiction: the spatial check in Definition~\ref{def:auth} is evaluated
against the jurisdiction $\ell$ declared in the request, never against the
jurisdiction of an intermediate junior role. This is deliberate and is the point
of Remark~\ref{rem:launder}.

\begin{definition}[Temporal authorisation graph]
The \emph{temporal authorisation graph} of $\Theta$ is the directed graph
$G = (V, E, \tau)$ with $V = U \cup R \cup O$ and
$E = UA \cup H \cup PA$, each edge labelled by its validity. The
\emph{snapshot} at $t$ is $G(t) = (V, E_t)$ with
$E_t = \{ e \in E \mid \chi_e(t) = 1 \}$.
\end{definition}

Unlike the bipartite graph of \citet{kokash2026compliance}, $G$ is tripartite
with an intra-role layer, so paths of length greater than one exist and
reachability is a meaningful question. In the bipartite formulation every path
has length one and Theorem~4 of that sketch was vacuous.

\subsection{Constraint predicates}

\begin{definition}[Region compatibility]
\label{def:region}
A request from jurisdiction $\ell$ is \emph{region-compatible} with object $o$,
written $\RegionOK(\ell, o)$, iff
\[
  \underbrace{\forall \ell_g \in \rho_G(o) : (\ell, \ell_g) \in \Gamma}
             _{\text{every governing regime permits}}
  \;\land\;
  \underbrace{\exists \ell_s \in \rho_S(o) : (\ell, \ell_s) \in \Gamma}
             _{\text{some replica is reachable}} .
\]
\end{definition}

\begin{remark}[Forum shopping under a single existential]
\label{rem:forum}
\citet{kokash2026compliance} define region compatibility as
$\exists \ell_r \in \rho_R(r), \exists \ell_o \in \rho_O(o) : (\ell_r, \ell_o)
\in \Gamma$, with $\rho_O$ described as the jurisdictions of the object's
storage locations but used as though it carried the governing regime. Under
that definition, an object replicated in an adequate and a non-adequate
jurisdiction is accessible via the adequate replica even when the non-adequate
regime governs it, and a multinational participant is accessible from whichever
of its offices happens to be adequate. Definition~\ref{def:region} closes both
by splitting the two roles of $\rho_O$ and by quantifying universally over
governance. When $|\rho_G(o)| = 1$ and $\rho_G(o) \subseteq \rho_S(o)$, the two
definitions coincide, so instances with a single governing regime per object,
including the running example of that paper, are unaffected.
\end{remark}

\begin{definition}[Purpose limitation and minimisation]
For $(r,o) \in PA^{*}_t$, a request satisfies \emph{purpose limitation} iff
$\PurpOK_t(r,o,p) \iff p \in \Pi_t(r,o)$. A requested object set
$O' \subseteq O$ satisfies \emph{data minimisation} for purpose $p$ iff
$\MinData(O',p) \iff O' \subseteq \Delta(p)$; for a single object this is
$o \in \Delta(p)$.
\end{definition}

\subsection{Authorisation and compliance}

\begin{definition}[STP-RBAC authorisation]
\label{def:auth}
A request $q = (u,r,o,p,\ell,t)$ is \emph{authorised} under $\Theta$, written
$\Auth_\Theta(q)$, iff all six conditions hold:
\[
\begin{aligned}
  \Auth_\Theta(u,r,o,p,\ell,t) \iff\;
  & \underbrace{r \in \Act_t(u)}_{\text{activation}}
  \;\land\;
    \underbrace{\ell \in \rho_U(u)}_{\text{locus}} \\[4pt]
  &\land\;
    \underbrace{(r,o) \in PA^{*}_t}_{\text{effective assignment}}
  \;\land\;
    \underbrace{\RegionOK(\ell,o)}_{\text{spatial}} \\[4pt]
  &\land\;
    \underbrace{\PurpOK_t(r,o,p)}_{\text{purpose}}
  \;\land\;
    \underbrace{o \in \Delta(p)}_{\text{minimisation}} .
\end{aligned}
\]
\end{definition}

Temporal validity does not appear as a separate conjunct because it is carried
inside $\Act_t$ and $PA^{*}_t$. The structural condition $(r,o) \in PA$ of the
earlier formulation is likewise subsumed: $PA^{*}_t$ is empty on pairs with no
assignment, so unassigned pairs are rejected rather than treated as merely
inactive.

\begin{proposition}[Conservative extension]
\label{prop:conservative}
Let $\Theta$ satisfy $|U| = 1$, $H = \emptyset$, $UA = \{u_0\} \times R$ with
$\tau_U \equiv \Tdom$, $|\rho_U(u_0)| = 1$, and
$\rho_G(o) = \rho_S(o)$ with $|\rho_G(o)| = 1$ for every $o$. Then for the
unique admissible $\ell$,
\[
  \Auth_\Theta(u_0, r, o, p, \ell, t)
  \iff
  \Auth^{\mathrm{prev}}(r,o,p,t),
\]
where $\Auth^{\mathrm{prev}}$ is the authorisation predicate of
\citet{kokash2026compliance}.
\end{proposition}

\begin{proof}
Under these restrictions $\Act_t(u_0) = R$ and the locus conjunct is vacuous.
With $H = \emptyset$ we have $\succeq_t$ = identity, so
$PA^{*}_t = \{(r,o) \in PA \mid \chi_{(r,o)}(t) = 1\}$ and
$\Pi_t = \Pi$, which reproduces the assignment and temporal conjuncts. With
$\rho_G(o) = \rho_S(o)$ a singleton, the universal and existential clauses of
Definition~\ref{def:region} collapse to the same literal, which is the earlier
$\RegionOK$. The purpose and minimisation conjuncts are unchanged.
\end{proof}

Proposition~\ref{prop:conservative} matters because it makes the present model
a refinement rather than a replacement: policies already written against the
earlier definitions retain their meaning.

\begin{definition}[Spatio-temporal privacy compliance]
$\Theta$ is \emph{compliant} with respect to a legality oracle $\LegalOK$ iff
$\Auth_\Theta(q) \Rightarrow \LegalOK(q)$ for every request $q$. A
\emph{violation} is a request $q$ with $\Auth_\Theta(q)$ and
$\neg\LegalOK(q)$.
\end{definition}

$\LegalOK$ is external to the model. It represents the ground truth that
regulation, ethics approval, and contract jointly determine, and it is what the
elicitation pipeline of \citet{kokash2026compliance} approximates. The model's
job is to make $\Auth_\Theta$ small enough to sit inside it, and the evaluation
in Section~\ref{sec:evaluation} measures exactly that containment.

% ============================================================
\section{Structural Results}
\label{sec:results}

\subsection{Snapshot reduction}

\begin{definition}[Projected snapshot]
For $t \in \Tdom$, $\ell \in \Ldom$, $p \in \Pdom$, the \emph{projected
snapshot} is the bipartite-layered graph
$G(t,\ell,p) = (U \cup R \cup O,\; A_t \cup P_t^{\ell,p})$ with
\begin{align*}
  A_t &= \{(u,r) \mid r \in \Act_t(u)\}, \\
  P_t^{\ell,p} &= \bigl\{ (r,o) \in PA^{*}_t \;\big|\;
      \RegionOK(\ell,o) \land p \in \Pi_t(r,o) \land o \in \Delta(p) \bigr\}.
\end{align*}
\end{definition}

\begin{theorem}[Snapshot reduction]
\label{thm:snapshot}
For every $u, r, o$ and every fixed $(t,\ell,p)$ with $\ell \in \rho_U(u)$,
\[
  \Auth_\Theta(u,r,o,p,\ell,t)
  \iff
  (u,r) \in A_t \;\land\; (r,o) \in P_t^{\ell,p},
\]
that is, iff $u$ reaches $o$ by a path of length two in $G(t,\ell,p)$.
Consequently $G(t,\ell,p)$ is a classical RBAC instance in the sense of
\citet{sandhu1996rbac}, with user assignment $A_t$ and permission assignment
$P_t^{\ell,p}$, and every decision procedure for flat RBAC applies to it
unchanged.
\end{theorem}

\begin{proof}
Immediate by unfolding Definition~\ref{def:auth}: the first conjunct is
membership in $A_t$, the locus conjunct holds by assumption, and the remaining
four conjuncts are exactly the defining conditions of $P_t^{\ell,p}$. Length-two
paths in a graph whose only edges run $U \to R$ and $R \to O$ are precisely
pairs of such edges sharing their middle vertex.
\end{proof}

This is the statement that Theorem~1 of the earlier sketch intended. That
version wrote $G(t,s,p)$ with an undefined parameter $s$ and asserted
equivalence to classical RBAC without specifying the projection, so it could not
be proved as stated.

\begin{corollary}[Soundness and completeness of two-phase enforcement]
\label{cor:split}
Define the coordinator prefilter
$G^{\mathrm{ST}}(t,\ell) = \{(r,o) \in PA^{*}_t \mid \RegionOK(\ell,o)\}$,
computed without knowledge of $p$, and let the worker node accept
$(r,o) \in G^{\mathrm{ST}}(t,\ell)$ iff
$p \in \Pi_t(r,o) \land o \in \Delta(p)$. The composite decision equals
$\Auth_\Theta$ exactly: no authorised request is filtered out, and no
unauthorised request survives both phases.
\end{corollary}

\begin{proof}
$\Auth_\Theta$ is a conjunction, and the two phases partition its conjuncts
into those independent of $p$ and those dependent on $p$. Conjunction is
associative, so the split preserves the value.
\end{proof}

Corollary~\ref{cor:split} is the formal justification for the deployment
architecture of \citet{kokash2026compliance}, in which temporal and spatial
gating happens centrally and purpose gating happens at the data provider. It
also delimits what the coordinator learns: it evaluates a purpose-independent
predicate, so it need not be told why a workflow step wants an object.

\subsection{Monotonicity and downward closure of compliance}

\begin{definition}[Relaxation order]
\label{def:relax}
For policies $\Theta, \Theta'$ over the same $U, R, O, \Ldom, \Pdom$, write
$\Theta \sqsubseteq \Theta'$ ($\Theta'$ \emph{relaxes} $\Theta$) iff
\begin{enumerate}
  \item[(i)] $UA \subseteq UA'$, \; $PA \subseteq PA'$, \; $H \subseteq H'$;
  \item[(ii)] $\tau_U(x) \subseteq \tau'_U(x)$, \; $\tau_P(x) \subseteq
        \tau'_P(x)$, \; $\tau_H(x) \subseteq \tau'_H(x)$, pointwise on the
        smaller domain;
  \item[(iii)] $\Gamma \subseteq \Gamma'$, \; $\rho_U(u) \subseteq \rho'_U(u)$,
        \; $\Pi(r,o) \subseteq \Pi'(r,o)$, \; $\Delta(p) \subseteq \Delta'(p)$;
  \item[(iv)] $\rho'_G(o) \subseteq \rho_G(o)$ \; and \;
        $\rho_S(o) \subseteq \rho'_S(o)$.
\end{enumerate}
\end{definition}

Note the reversal on $\rho_G$. Governing regimes appear under a universal
quantifier, so \emph{removing} one weakens the gate; every other component
appears positively and is relaxed by enlargement. This asymmetry is easy to get
wrong and is the reason the informal claim ``extending validity intervals,
purpose scope or spatial area never removes permissions'' in the earlier draft
is not correct as stated: enlarging $\rho_G$ removes permissions.

\begin{theorem}[Monotonicity]
\label{thm:monotone}
If $\Theta \sqsubseteq \Theta'$ then
$\Auth_\Theta(q) \Rightarrow \Auth_{\Theta'}(q)$ for every request $q$.
\end{theorem}

\begin{proof}
We show each conjunct of Definition~\ref{def:auth} is monotone in
$\sqsubseteq$.

\emph{Activation.} $UA \subseteq UA'$ and $\tau_U \subseteq \tau'_U$ give
$\chi^{\Theta}_{(u,r')}(t) \le \chi^{\Theta'}_{(u,r')}(t)$. Since
$H \subseteq H'$ and $\tau_H \subseteq \tau'_H$, we get
$H_t \subseteq H'_t$ and hence ${\succeq_t} \subseteq {\succeq'_t}$, because the
reflexive-transitive closure is monotone in its argument. Therefore
$\Act_t(u) \subseteq \Act'_t(u)$.

\emph{Locus.} $\rho_U(u) \subseteq \rho'_U(u)$ directly.

\emph{Effective assignment and purposes.} $PA \subseteq PA'$,
$\tau_P \subseteq \tau'_P$, and ${\succeq_t} \subseteq {\succeq'_t}$ give
$PA^{*}_t \subseteq PA'^{*}_t$; the union defining $\Pi_t$ is over a larger
index set of larger sets, so $\Pi_t(r,o) \subseteq \Pi'_t(r,o)$.

\emph{Spatial.} Suppose $\RegionOK^{\Theta}(\ell,o)$. For the universal clause,
$\rho'_G(o) \subseteq \rho_G(o)$ means the clause ranges over fewer labels, and
$\Gamma \subseteq \Gamma'$ means each surviving obligation is still met. For the
existential clause, the witness $\ell_s \in \rho_S(o)$ lies in
$\rho'_S(o)$ and $(\ell,\ell_s) \in \Gamma \subseteq \Gamma'$.

\emph{Purpose and minimisation.} Immediate from
$\Pi_t \subseteq \Pi'_t$ and $\Delta(p) \subseteq \Delta'(p)$.
\end{proof}

\begin{corollary}[Compliance is downward closed]
\label{cor:downward}
If $\Theta \sqsubseteq \Theta'$ and $\Theta'$ is compliant, then $\Theta$ is
compliant.
\end{corollary}

\begin{proof}
Contrapositive of Theorem~\ref{thm:monotone}: if $q$ violates $\Theta$ then
$\Auth_\Theta(q)$, hence $\Auth_{\Theta'}(q)$, and $\neg\LegalOK(q)$ is
independent of the policy, so $q$ violates $\Theta'$.
\end{proof}

Corollary~\ref{cor:downward} has an operational reading that we exploit in
Section~\ref{sec:evaluation}. Verify once against the most permissive envelope
a project will ever adopt. Thereafter, any change that only tightens the policy,
which covers most onboarding events, needs no re-verification. Only relaxations
trigger a re-check, and the check can be restricted to the difference. This
turns full re-verification on every participant join into an incremental
operation.

\subsection{Temporal separation of duty}

The earlier draft defined a conflict predicate over
$\Auth(r_1,o_1,t) \land \Auth(r_2,o_2,t)$ and concluded that a conflict exists
iff the corresponding validity intervals intersect. That statement is not
well-formed: separation of duty constrains what a single principal may hold, and
with no users in the signature there is nothing to bind $r_1$ and $r_2$
together. Two different people holding conflicting roles at the same time is
normal operation, not a violation. The corrected statement follows.

\begin{definition}[Temporal SoD violation]
$\Theta$ admits a \emph{temporal SoD violation} iff there exist $u \in U$,
$(r_1,r_2) \in \SoD$, and $t \in \Tdom$ with
$r_1, r_2 \in \Act_t(u)$.
\end{definition}

\begin{lemma}[Piecewise constancy]
\label{lem:breakpoints}
Let $B \subset \Tdom$ be the set of endpoints of the constituent intervals of
all $\tau_U$ and $\tau_H$ labels, and let $\mathcal{A}$ be the arrangement of
$\Tdom$ induced by $B$, that is, the partition into maximal subintervals
containing no point of $B$ in their interior. Then $t \mapsto \Act_t(u)$ is
constant on each cell of $\mathcal{A}$, for every $u$.
\end{lemma}

\begin{proof}
$\Act_t(u)$ depends on $t$ only through the family of indicators
$\{\chi_{(u,r')}(t)\}$ and $\{\chi_{(r,r')}(t)\}$, since $\succeq_t$ is a
function of $H_t$ alone. Each indicator is the characteristic function of a
finite union of closed intervals with endpoints in $B$, hence constant on the
interior of every cell of $\mathcal{A}$ and on every degenerate cell
$\{b\}$, $b \in B$. A constant family yields a constant closure.
\end{proof}

\begin{theorem}[SoD reduces to interval stabbing]
\label{thm:sod}
Let $N = |UA| + |H|$ counted with interval multiplicity. A temporal SoD
violation exists iff one exists at some point of a set of at most $2N$
candidate instants, namely one representative per cell of $\mathcal{A}$. In the
hierarchy-free case $H = \emptyset$, a violation exists iff there are
$u$ and $(r_1,r_2) \in \SoD$ with
$\tau_U(u,r_1) \cap \tau_U(u,r_2) \neq \emptyset$, which for constituent
intervals $[a_1,b_1]$ and $[a_2,b_2]$ holds iff
$\max(a_1,a_2) \le \min(b_1,b_2)$.
\end{theorem}

\begin{proof}
By Lemma~\ref{lem:breakpoints}, $\Act_t(u)$ takes at most $2N$ distinct values
as $t$ ranges over $\Tdom$, one per cell, so it suffices to test one
representative per cell; $|\mathcal{A}| \le 2N+1$ and the unbounded outer cells
have empty activation. For $H = \emptyset$ we have
$\Act_t(u) = \{r \mid (u,r) \in UA, \chi_{(u,r)}(t) = 1\}$, so
$r_1, r_2 \in \Act_t(u)$ for some $t$ iff the two validity sets share a point,
and two closed intervals share a point exactly under the stated inequality.
\end{proof}

Theorem~\ref{thm:sod} is what makes SoD checking a sweep-line problem rather
than a search over $\Tdom$: sort the $2N$ endpoints, maintain the active set
incrementally, and test the $\SoD$ pairs against it. Section~\ref{sec:algorithms}
gives the bound.

\subsection{Workflow feasibility}

Definition~\ref{def:auth} answers a question about one instant. Requirement R5
asks a different question: given a workflow whose steps have their own role,
object, purpose, and location demands, and a precedence order among them, can
the workflow run at all? The earlier draft invoked time-respecting paths in the
sense of \citet{holme2012temporal} for the single-access case, which is a
category error. A single authorisation needs all its edges live
\emph{simultaneously}, which is Theorem~\ref{thm:snapshot}, a statement about a
snapshot. Time-respecting paths are the right notion for workflows, where a
later step may rely on a permission that was not yet active when an earlier step
ran.

\begin{definition}[Workflow]
A \emph{workflow} is $W = (S, \prec, \mathit{dem})$ where $S$ is a finite set of
steps, $\prec$ is a strict partial order on $S$ given as a DAG, and
$\mathit{dem}(s) = (u_s, r_s, o_s, p_s, \ell_s)$ gives the demand of step $s$. A
\emph{schedule} is $\sigma : S \to \Tdom$ with $s \prec s' \Rightarrow
\sigma(s) < \sigma(s')$. $W$ is \emph{feasible} under $\Theta$ iff some schedule
$\sigma$ satisfies $\Auth_\Theta(u_s, r_s, o_s, p_s, \ell_s, \sigma(s))$ for
every $s \in S$.
\end{definition}

\begin{theorem}[Workflow feasibility]
\label{thm:workflow}
Let $F_s = \{ t \mid \Auth_\Theta(u_s,r_s,o_s,p_s,\ell_s,t) \} \in \mathcal{I}$
be the feasible window of step $s$, a finite union of at most $m$ intervals.
Then:
\begin{enumerate}
  \item $W$ is feasible iff there is a time-respecting path through the
        step-expanded graph, and feasibility is decidable in
        $O\bigl((|S| + |{\prec}|)\, m \log m\bigr)$ time by earliest-start
        propagation in topological order, which additionally returns the
        pointwise-earliest feasible schedule.
  \item If step demands specify a \emph{set} of admissible roles per step rather
        than a single role, and $\SoD$ constraints are imposed across steps,
        deciding feasibility is NP-complete.
\end{enumerate}
\end{theorem}

\begin{proof}
(1) $F_s$ is computable from Definition~\ref{def:auth} because each conjunct is
either time-independent or a finite union of intervals, and finite unions are
closed under intersection. Process $S$ in topological order maintaining
$e(s) = $ earliest feasible start. Set
$e(s) = \min \{ t \in F_s \mid t > \max_{s' \prec s} e(s') \}$, computed by
binary search in the sorted interval list of $F_s$, and report infeasible if no
such $t$ exists. Earliest-start propagation on a DAG is optimal for this
constraint shape by the standard exchange argument: if any feasible schedule
exists, the pointwise-earliest one is feasible, since delaying a predecessor can
only shrink the admissible set of a successor. Each step performs one binary
search over $m$ intervals, and each precedence edge is relaxed once.

(2) Membership in NP is clear: a schedule together with a role assignment is a
polynomial certificate checkable by (1). Hardness follows by reduction from the
workflow satisfiability problem, which \citet{wang2010wsp} show is NP-complete
even without a temporal dimension. Given a WSP instance, take $\Tdom$ a single
instant, $\Gamma$ total, $\Pi$ and $\Delta$ unconstrained, all validity windows
equal to $\Tdom$, and $\prec$ empty; then feasibility of $W$ holds exactly when
the WSP instance is satisfiable. The construction is polynomial.
\end{proof}

Part (2) explains why the earlier draft's complexity table could not be right.
The table listed static consistency as ``NP-complete (interval arithmetic is
polytime, but with hierarchy becomes NP-hard)'', conflating two different
problems and two different sources of hardness. Hardness does not come from the
hierarchy, which only adds a transitive closure. It comes from choice: as long
as every step's principal and role are fixed, everything is polynomial. As soon
as the policy must be searched for an assignment, the problem inherits WSP
hardness. Fixed-parameter results for WSP in the number of
steps~\citep{crampton2013wspseniority} carry over and are the practical route.

% ============================================================
\section{Algorithms and Complexity}
\label{sec:algorithms}

Let $n = |U| + |R| + |O|$, let $M = |UA| + |PA| + |H|$ counted with interval
multiplicity, and let $k = |\SoD|$. We assume each edge's validity is stored as
a sorted list of disjoint intervals, and that an interval tree over all validity
intervals is built once in $O(M \log M)$.

\begin{description}
  \item[A1: Authorisation query.] Given $q$, decide $\Auth_\Theta(q)$. With the
        interval tree and a precomputed $\succeq_t$ for the current cell of
        $\mathcal{A}$, this is $O(\log M + d)$ where $d$ is the number of roles
        junior to $r$ that hold $o$. Without precomputation it is
        $O(|H| + \log M)$, dominated by the closure walk. The earlier draft's
        $O(\log n)$ holds only for the hierarchy-free case with a single
        candidate edge.
  \item[A2: SoD conflict detection.] Report all temporal SoD violations. Sort
        the $2M$ endpoints, sweep, and maintain $\Act_t(u)$ incrementally per
        user. Testing $\SoD$ pairs against the active set at each of the $O(M)$
        breakpoints gives $O(M \log M + Mk)$ in the worst case, and
        $O(M \log M + Z)$ output-sensitive when $\SoD$ is indexed by role, where
        $Z$ is the number of reported violations. The earlier draft's nested
        loop over $PA \times PA$ is $O(|PA|^2)$ and also tests the wrong
        condition, since it does not quantify over users.
  \item[A3: Workflow feasibility, fixed roles.] Theorem~\ref{thm:workflow}(1),
        $O((|S| + |{\prec}|) m \log m)$.
  \item[A4: Policy consistency.] Check well-formedness: every validity set is
        non-empty and each constituent interval satisfies $t^s \le t^e$; $H$ is
        acyclic; $\Pi(r,o) \neq \emptyset$ for every $(r,o) \in PA$; and
        $\Pi(r,o) \subseteq \{p \mid o \in \Delta(p)\}$, that is, no permission
        grants a purpose that minimisation will always deny. This last check is
        worth stating separately because such a pair is dead policy: it can
        never authorise anything and almost always indicates an authoring error.
        All checks are $O(M + |PA| \cdot |\Pdom|)$.
  \item[A5: Reachability envelope.] Compute, for a given $(\ell,p)$, the set of
        objects any user can reach over a horizon. Union the projected snapshots
        over the cells of $\mathcal{A}$, $O(M \log M + M n)$ in the worst case.
\end{description}

Table~\ref{tab:complexity} summarises. The essential distinction, absent from
the earlier draft, is between \emph{checking} a fully instantiated policy, which
is polynomial throughout, and \emph{searching} for interval assignments or role
assignments, which is not.

\begin{table}[t]
\centering
\caption{Complexity of the analysis problems. $M$ is the number of assignment
edges counted with interval multiplicity, $Z$ the output size. The last row is
not our result and is quoted from \citet{sasturkar2011arbac}.}
\label{tab:complexity}
\small
\begin{tabular}{llll}
\toprule
\textbf{Problem} & \textbf{Input} & \textbf{Complexity} & \textbf{Method} \\
\midrule
Authorisation query   & ground policy & $O(\log M + d)$        & interval tree \\
SoD detection         & ground policy & $O(M\log M + Z)$       & sweep line \\
Workflow feasibility  & fixed roles   & $O((|S|+|{\prec}|)m\log m)$ & DAG propagation \\
Workflow feasibility  & role choice   & NP-complete            & SMT / FPT in $|S|$ \\
Policy consistency    & ground policy & $O(M + |PA||\Pdom|)$   & direct \\
Interval synthesis    & free windows  & NP-hard                & SMT \\
Admin.\ reachability  & ARBAC         & PSPACE-complete        & symbolic \\
\bottomrule
\end{tabular}
\end{table}

% ============================================================
\section{SMT Encoding}
\label{sec:smt}

We give an encoding into quantifier-free linear integer arithmetic with
uninterpreted predicates, discharged by Z3~\citep{demoura2008z3}. The earlier
draft encoded only the temporal dimension and left the spatial and purpose
dimensions as a TODO; the encoding below covers all four. We assume $\Tdom$
discrete, which is standard for policy horizons measured in days or months.

\paragraph{Two modes.} The same encoding serves two problems, and keeping them
apart is what resolves the complexity confusion noted above.
\emph{Checking} fixes every interval endpoint as a numeral and asks whether a
property holds; the resulting formula is ground and decided in polynomial time.
\emph{Synthesis and repair} leaves endpoints as free integer variables and asks
for a model; this is where NP-hardness enters, since forbidden-overlap
constraints are negated conjunctions of linear inequalities and a set of them
encodes propositional satisfiability directly.

\paragraph{Variables.} For each $x \in UA \cup PA \cup H$ with constituent
interval index $j$, integers $a_{x,j}, b_{x,j}$ with the well-formedness
constraint $a_{x,j} \le b_{x,j}$. Booleans $\Gamma_{\ell,\ell'}$ for the
jurisdiction relation, $\Pi_{r,o,p}$ for purposes, $\Delta_{p,o}$ for
minimisation, $\rho^{G}_{o,\ell}$ and $\rho^{S}_{o,\ell}$ for governance and
storage, and $\rho^{U}_{u,\ell}$ for the user locus. In checking mode all
Booleans are fixed by the policy.

\paragraph{Activation.} $\mathit{live}_x(t) \leftrightarrow
\bigvee_j (a_{x,j} \le t \land t \le b_{x,j})$.

\paragraph{Hierarchy closure.} Introduce $\mathit{Sen}_{r,r'}(t)$ with the
Horn rules
\[
  \mathit{Sen}_{r,r}(t), \qquad
  \mathit{Sen}_{r,r''}(t) \leftarrow
    \mathit{Sen}_{r,r'}(t) \land \mathit{live}_{(r',r'')}(t),
\]
which the fixed-point engine of Z3 discharges directly. For a bounded role
depth $D$, the rules may instead be unrolled $D$ times into quantifier-free
form, which is preferable when the rest of the query is ground.

\paragraph{Spatial predicate.}
\[
  \RegionOK_{o,\ell}
  \;\leftrightarrow\;
  \Bigl(\bigwedge_{\ell_g \in \Ldom}
      \bigl(\rho^{G}_{o,\ell_g} \rightarrow \Gamma_{\ell,\ell_g}\bigr)\Bigr)
  \;\land\;
  \Bigl(\bigvee_{\ell_s \in \Ldom}
      \bigl(\rho^{S}_{o,\ell_s} \land \Gamma_{\ell,\ell_s}\bigr)\Bigr).
\]

\paragraph{Authorisation.}
\[
\begin{aligned}
  \Auth_{u,r,o,p,\ell}(t) \;\leftrightarrow\;
  & \Bigl(\bigvee_{r' \in R}
      \mathit{live}_{(u,r')}(t) \land \mathit{Sen}_{r',r}(t)\Bigr)
    \;\land\; \rho^{U}_{u,\ell} \\[4pt]
  &\land\;
    \Bigl(\bigvee_{r'' \in R}
      \mathit{Sen}_{r,r''}(t) \land \mathit{live}_{(r'',o)}(t)
      \land \Pi_{r'',o,p}\Bigr) \\[4pt]
  &\land\; \RegionOK_{o,\ell} \;\land\; \Delta_{p,o}.
\end{aligned}
\]
The third conjunct fuses effective assignment with effective purpose, since both
quantify over the same junior role $r''$. Splitting them, as a naive
transcription of Definition~\ref{def:effective} would, is unsound: it would
allow the assignment witness and the purpose witness to be different roles.

\paragraph{Queries.} Each analysis problem becomes a satisfiability query.
\begin{itemize}
  \item \emph{Unauthorised access}: assert $\Auth_{u,r,o,p,\ell}(t)$ for a
        request the policy author believes forbidden; a model is a
        counterexample.
  \item \emph{Temporal SoD}: assert
        $\exists t: \bigvee_{r'} \mathit{live}_{(u,r')}(t) \land
        \mathit{Sen}_{r',r_1}(t)$ conjoined with the same for $r_2$, over
        $(r_1,r_2) \in \SoD$.
  \item \emph{Dead permission}: assert
        $\neg \exists t,\ell : \Auth_{u,r,o,p,\ell}(t)$ for each
        $(r,o) \in PA$ and $p \in \Pi(r,o)$; unsatisfiability of the negation
        certifies the permission is reachable.
  \item \emph{Repair}: fix the Boolean policy structure, free the interval
        endpoints, assert the desired safety properties plus the availability
        properties the project needs, and minimise total window shrinkage as an
        optimisation objective.
\end{itemize}

\begin{remark}
Purpose limitation and minimisation are the only dimensions expressible without
arithmetic. A deployment that needs only these should use a
policy-as-code engine such as OPA/Rego~\citep{opa} rather than an SMT solver,
which is the pragmatic reading of Table~\ref{tab:verification}: the solver earns
its cost only when interval reasoning or synthesis is involved.
\end{remark}

% ============================================================
\section{Verification Architecture}
\label{sec:architecture}

Table~\ref{tab:verification} maps model constructs to tooling. Each row states
what the component decides and what it cannot decide, since a mapping that
lists technologies without their scope invites the reader to assume more
coverage than exists.

\begin{table}[t]
\centering
\caption{Verification components and their scope.}
\label{tab:verification}
\small
\begin{tabular}{p{2.4cm}p{5.2cm}p{4.6cm}}
\toprule
\textbf{Component} & \textbf{Decides} & \textbf{Technology and scope} \\
\midrule
Policy model &
Serialisation of $\Theta$; well-formedness A4 &
JSON Schema plus graph store; syntactic only \\
\addlinespace
Constraint engine &
Interval satisfiability, SoD, dead permissions, repair &
Z3~\citep{demoura2008z3}; ground checking polynomial, synthesis NP-hard \\
\addlinespace
Bounded consistency &
Structural counterexamples over small scopes &
Alloy~\citep{jackson2002alloy}; small-scope only, no proof of absence \\
\addlinespace
Temporal properties &
Safety and liveness over policy evolution &
NuSMV~\citep{cimatti2002nusmv} or TLA$^{+}$~\citep{lamport2002tla}; requires a finite abstraction of $\Tdom$ \\
\addlinespace
Runtime enforcement &
Per-request $\Auth_\Theta$ at coordinator and worker &
OPA/Rego~\citep{opa} and the Brane policy reasoner~\citep{kokash2024reasoner} \\
\addlinespace
Explanation &
Which conjunct failed and which policy element caused it &
Unsat cores and Z3 proofs, rendered to text \\
\bottomrule
\end{tabular}
\end{table}

The explanation layer deserves comment because it is where a formal model earns
its keep for non-specialist users. Definition~\ref{def:auth} is a conjunction of
six named conjuncts, so a denial always has an identifiable cause, and
Corollary~\ref{cor:split} tells us which node observed it. A denial can
therefore be reported as ``the request was refused at the provider because the
declared purpose is not among those permitted for this role on this object'',
with the offending $\Pi(r,o)$ entry attached. We treat natural-language
rendering of this structured diagnosis as a presentation concern, not a
reasoning one: the diagnosis itself comes from the unsat core, not from a
language model.

% ============================================================
\section{Evaluation Design}
\label{sec:evaluation}

The claims that require empirical support are: that the model expresses the
constraints real FDP projects impose (expressiveness); that checking is fast
enough for the coordinator to sit on the request path (performance); that the
incremental strategy licensed by Corollary~\ref{cor:downward} actually pays
(incrementality); and that the diagnostics are usable by the people who author
policies (usability). We evaluate each against a stated baseline.

\paragraph{E1: Expressiveness.} Encode the pharmaceutical and hospital
collaboration scenario of \citet{kokash2026compliance}, which has five roles,
six object classes, five purposes, four project phases, and five jurisdictions,
together with two further scenarios drawn from published federated healthcare
deployments. Report the fraction of stated constraints expressible in $\Theta$
without extension, and characterise the residue. We expect obligations and
consent revocation to fall outside, since STP-RBAC is a permission model and not
a normative one in the sense of eFLINT~\citep{binsbergen2020eflint}.

\paragraph{E2: Checking performance.} Measure A1, A2, and A4 on synthetic
policies scaling $|R|$, $|O|$, and interval multiplicity independently, against
a naive baseline that evaluates $\Auth$ by linear scan. Report median and
tail latency for A1, since coordinator admission is on the critical path.

\paragraph{E3: Solver scaling.} Measure Z3 runtime for the four query classes of
Section~\ref{sec:smt} as policy size grows, separately in checking and synthesis
mode. The hypothesis under test is that the two modes have visibly different
scaling curves, which is the empirical counterpart of the complexity
distinction.

\paragraph{E4: Incrementality.} Simulate participant onboarding sequences and
compare full re-verification against the incremental strategy: classify each
policy delta as a tightening or a relaxation under Definition~\ref{def:relax},
skip verification on tightenings, and verify only the difference on relaxations.
Report the fraction of onboarding events that are pure tightenings and the
resulting reduction in solver invocations. This experiment directly measures the
value of Corollary~\ref{cor:downward}, and it fails informatively: if most real
deltas turn out to be relaxations, the corollary is sound but operationally
uninteresting.

\paragraph{E5: Fault injection.} Seed policies with the four defect classes the
model is meant to catch (expired window still granting access, cross-border
grant with no $\Gamma$ edge, purpose creep, minimisation violation) and measure
detection rate and time to detection. Include the two defects the corrected
model catches but the earlier formulation did not, namely the forum-shopping
pattern of Remark~\ref{rem:forum} and the hierarchy-laundering pattern of
Remark~\ref{rem:launder}, and confirm that the earlier formulation misses them.

\paragraph{E6: Diagnostic usability.} Present denial diagnoses to participants
in the two roles that matter, data stewards and workflow authors, and measure
whether they identify the responsible policy element and propose a correct fix.
Compare structured diagnoses against a raw solver trace.

Section~\ref{sec:conclusion} notes which of these are complete.

\begin{remark}[Hierarchy laundering]
\label{rem:launder}
Consider $r_{\mathrm{senior}} \succeq_t r_{\mathrm{junior}}$ where the junior
role's holders operate in an adequate jurisdiction and the senior role's holders
do not. Under a model that evaluates region compatibility against the role
whose permission is being inherited, the senior role acquires cross-border
access it could not have been granted directly. Definition~\ref{def:auth} avoids
this by evaluating $\RegionOK$ against the requester's declared $\ell$ alone,
independently of the inheritance path. This is why $\ell$ belongs to the request
rather than to the role.
\end{remark}

% ============================================================
\section{Related Work}
\label{sec:related}

\paragraph{Temporal RBAC.} \citet{bertino2001trbac} introduced TRBAC,
associating time constraints with role enabling and permission assignment, with
periodic schedules, calendar expressions, and temporal separation of duty. The
enabled role-permission relation at time $t$ is the direct ancestor of our
$\chi_e(t)$ and of the snapshot $G(t)$. Earlier work on temporal authorisation
bases~\citep{bertino2000tab} established interval-based authorisation
propagation. \citet{joshi2005gtrbac} generalised TRBAC to cover role enabling,
role activation, and permission assignment simultaneously, with triggered events
and a temporal-logic semantics. Our $\tau_U$, $\tau_P$, and $\tau_H$ correspond
to GTRBAC's three constraint families, and our restriction to finite unions of
intervals is the bounded-horizon unfolding of their periodic expressions.

\paragraph{Spatial and spatio-temporal RBAC.} GEO-RBAC~\citep{bertino2005georbac}
makes role activation depend on the user's position within spatial extents.
LoT-RBAC~\citep{chandran2005lotrbac} and the model of
\citet{ray2007strbac} combine location and time.
\citet{toahchoodee2009alloy} analyse a spatio-temporal model in Alloy. Our
spatial dimension differs in kind rather than degree: these models treat
location as physical position, whereas $\Gamma$ is a legal relation over
jurisdictions, asymmetric because adequacy decisions are, and attached to
objects in two distinct ways (governance and storage) because data localisation
and data governance are separate obligations. To our knowledge the
governance-versus-storage distinction has not been made in this literature.

\paragraph{Purpose-based access control.} \citet{byun2008purpose} associate
intended purposes with data elements and access purposes with requests,
supporting explicit prohibitions and purpose hierarchies. Our $\Pi$ is their
intended-purpose labelling lifted to role-object pairs, and $\Delta$ adds the
minimisation direction, which their model does not address. A purpose hierarchy
would be a natural extension of $\Pi$ and is left to future work.

\paragraph{Policy analysis and verification.} \citet{sasturkar2011arbac}
establish that user-role reachability under administrative RBAC is
PSPACE-complete, with tractable fragments under restrictions.
\citet{uzun2012trbacanalysis} and \citet{jha2014trbacadmin} extend safety
analysis to the temporal setting, and \citet{ranise2014atrbac} give scalable
symbolic procedures for administrative TRBAC.
\citet{mondal2009timedautomata} analyse GTRBAC via timed automata.
\citet{wang2010wsp} establish NP-completeness of the workflow satisfiability
problem, which Theorem~\ref{thm:workflow}(2) reduces from, and
\citet{crampton2013wspseniority} give fixed-parameter algorithms in the number
of steps. \citet{gouglidis2023modelchecking} apply model checking to a
production cloud IAM policy set. Our contribution here is not new complexity
theory but the placement of our specific analysis problems relative to these
known boundaries, which the earlier draft did not attempt.

\paragraph{Temporal graphs.} Our $G(t)$ is a temporal
network~\citep{holme2012temporal}, and time-respecting paths are the right
formal notion for workflow feasibility. We stress the boundary established in
Section~\ref{sec:results}: instantaneous authorisation is a snapshot question
and does not need this machinery, while workflow feasibility does.

\paragraph{Federated data processing.} Brane~\citep{valkering2021brane} provides
the execution substrate, with policy enforcement via
eFLINT~\citep{binsbergen2020eflint} and OPA/Rego~\citep{opa,kokash2024reasoner}.
JustAct~\citep{esterhuyse2024justact} addresses justification and audit of
actions against inter-organisational policies, which is complementary: it asks
whether an action can be justified, while we ask whether a policy admits actions
it should not. The EPI project~\citep{kassem2024epi} supplies the healthcare
deployment context, and \citet{kokash2026compliance} the compliance-management
architecture this model formalises.

% ============================================================
\section{Discussion and Conclusion}
\label{sec:conclusion}

We have given STP-RBAC a signature that supports the questions the FDP setting
actually raises, and proved the four structural results the model needs to be
usable: a reduction to classical RBAC per snapshot, monotonicity with its
incremental-verification corollary, an interval characterisation of temporal
separation of duty, and a placement of workflow feasibility on the tractable
side of the boundary when roles are fixed and the intractable side when they are
not.

Several limitations are worth stating plainly. The model is a permission model:
it says what may happen, not what must. Obligations, consent revocation with
retroactive effect, and rights of erasure are normative constructs better served
by eFLINT~\citep{binsbergen2020eflint}, and a hybrid deployment is the realistic
target. The legality oracle $\LegalOK$ is an idealisation, and the gap between
it and any implementable approximation is exactly where the compliance risk
lives; our model can prove $\Auth_\Theta$ sound against a stated $\LegalOK$, but
it cannot validate $\LegalOK$ itself. $\Gamma$ is a flat relation over labels
and does not model conditional adequacy, standard contractual clauses, or
derogations, all of which are conditional rather than binary in practice.
Finally, the evaluation of Section~\ref{sec:evaluation} is a design; E1 and E2
are in progress against the BraneHub prototype~\citep{kokash2026branehub}, and
E3 through E6 are not yet run.

Future work runs in three directions: purpose hierarchies with a subsumption
order on $\Pdom$, which would let minimisation be stated once per purpose family
rather than per purpose; conditional jurisdiction, replacing $\Gamma$ with
guarded edges carrying safeguard requirements; and integration of the workflow
feasibility check into the Brane orchestration path, so that a workflow is
rejected at submission with a specific reason rather than failing at the first
denied step.

\bibliographystyle{plainnat}
\bibliography{main}

\end{document}
